\section{Introduction}\label{sec:introduction}

Stationarity is the weakest variational condition naturally imposed on a
minimal submanifold: the first variation of area vanishes under every
compactly supported ambient deformation. Integral varifolds retain this
condition under weak limits while allowing integer multiplicity,
self-intersections, and singularities. Their general regularity problem
asks how much of the support must nevertheless be a smooth embedded
minimal submanifold.

Let $m,n\ge1$ and let $U\subset\R^{m+n}$ be open. An integral
$m$-varifold $V$ has weight
$\mu=\theta\Hh^m\llcorner E$, where $E$ is countably $m$-rectifiable
and $\theta$ is positive, integer-valued, and locally integrable.
At $\mu$-almost every $x$, write $T_xV$ for the approximate tangent
$m$-plane and $P_{T_xV}$ for its orthogonal projection. Stationarity means
\begin{equation}\label{eq:stationarity}
 \int\tr(P_{T_xV}DX(x))\,\dd\mu(x)=0
 \qquad\text{for every }X\in C_c^1(U;\R^{m+n}).
\end{equation}
A point $p\in\spt\mu\cap U$ is \emph{regular} if some neighborhood $W$
of $p$ contains a smooth embedded minimal $m$-submanifold $\Sigma$,
without boundary in $W$, such that $\spt\mu\cap W=\Sigma$ and
$V\llcorner W=q|\Sigma|$ for a positive integer $q$.
We denote the remaining support points by
$\Sing V$. Thus regularity includes equality of the local support with
one embedded submanifold, not merely a smooth parametrization on a set
of full measure.
For stationary integral varifolds this is equivalent to local $C^1$
embeddedness of the support, by constancy and elliptic regularity.

\begin{theorem}\label{thm:main}
For every stationary integral $m$-varifold $V$ in an open set
$U\subset\R^{m+n}$, where $m,n\ge1$,
\[
 \dim_{\Hh}\Sing V\le m-1.
\]
\end{theorem}

Theorem~\ref{thm:main} proves the Euclidean conjecture recorded as
Conjecture~1.1 by Brena, Decio, and De Lellis \cite{BDD2025}.
There is no stability, minimizing, orientability, density-upper-bound,
or codimension hypothesis. The estimate is sharp: the sum of the
varifolds of two distinct $m$-planes meeting along an $(m-1)$-plane is
stationary, and its singular set is exactly that intersection.

\subsection{Context and antecedents}
Allard's compactness and regularity theory \cite{Allard1972} supplies
the basic passage from weak convergence to geometric information and
gives a relatively open dense regular set. The dimension-reduction
method originating in \cite{Federer1970} leads to the classical
stratification by translation symmetries of tangent cones. For stationary
integral varifolds, the set of points admitting no planar tangent has
Hausdorff dimension at most $m-1$. Naber and Valtorta
\cite{NV2020} proved rectifiability of these classical strata for
integral varifolds with bounded mean curvature. The remaining difficulty
is therefore the singular set at higher-multiplicity planar tangents,
where Allard's multiplicity-one regularity theorem does not apply.
Standard treatments of the compactness, monotonicity, and tangent-cone
theory used here include \cite{Federer1969,Simon1983}.

For area-minimizing integral currents, Almgren's regularity theory
\cite{Almgren2000} gives an interior singular-set bound of dimension $m-2$ in
arbitrary codimension. De Lellis and Spadaro developed a new proof based
on multiple-valued functions, center manifolds, and a frequency analysis
of the blow-up \cite{DLS2011,DLS2016II,DLS2016III}.
These techniques explain the role of subtracting a well-chosen smooth
center before normalizing the height. Their regularity theorems do
not apply to arbitrary stationary integral varifolds; even the crossing
example above prevents the same dimension bound in our setting.

Hirsch and Spolaor proved the $m-1$ bound for two-valued stationary
Lipschitz graphs \cite{HSgraphs}. Their use of generalized-gradient
Young measures retains height and gradient information under weak
convergence and is a methodological antecedent of our measure-valued
blow-up, which does not assume graphicality or two-valuedness.
Hirsch and Spolaor also proved almost-everywhere regularity for stationary
codimension-one integral currents whose density is at most two at every support
point \cite{HS2025}. Krummel, Minter, and Wickramasekera
\cite{KMW2026} obtained singular-set bounds of dimension at most $m-1$ at
controlled densities in classes satisfying an explicit
epsilon-regularity assumption.

A different line of work addresses higher-order rectifiability.
Brakke's work gives $C^{1,\alpha}$ rectifiability, for every
$0<\alpha<1$, of stationary integral varifolds \cite{Brakke1978}.
Menne proved second-order rectifiability under locally bounded first
variation \cite{Menne2013}. Brena, De Lellis, and Franceschini
\cite{BDLF2025} established smooth rectifiability for stationary
integral varifolds; Kolasi\'nski \cite{Kolasinski2026} subsequently
gave an independent alternative proof, following a program he
attributes to Menne. Smooth rectifiability describes
the measure up to a null
set by smooth pieces; it does not assert that the entire support near
almost every point equals one such piece. This distinction is essential
for the singular set defined above.

The ambient geometry is also important. De Lellis, Hirsch, Lihn, and
Spolaor \cite{DLHLS2026} construct stationary integral varifolds with
planar tangents and infinite local topology in smooth, generally
non-Euclidean metrics. Their examples do not contradict the Euclidean
dimension bound proved here.

Brena, Decio, and De Lellis \cite[Conjectures~1.1 and~1.2]{BDD2025}
separate the dimension bound from the weaker assertion
$\Hh^m(\Sing V)=0$. The latter Euclidean assertion is proved in
\emph{Almost-everywhere regularity of stationary integral varifolds}
\cite{AE}, which we henceforth call AE. Our proof uses AE's signed-excess
theorem and local estimates from its fitting and frequency arguments.
The elliptic correction of the averaged graph in
\cite[Sections~6.3--6.4]{BDLF2025} is an antecedent of this fitting
construction. The almost-everywhere
regularity conclusion itself is not an input: a set can have zero
$\Hh^m$ measure and still have Hausdorff dimension greater than $m-1$.
We state the imported measure estimates in
Section~\ref{sec:normal-inputs}. Section~\ref{sec:fitting} specifies
the local fitting estimates and verifies the hypotheses needed for
their use at the prescribed scales.

\subsection{The proof and its new components}
Suppose that $\dim_{\Hh}\Sing V>m-1$ and choose an exponent $s$ strictly
between these two numbers. A first dimension-reduction argument produces
a singular point, placed at the origin, an integer $Q$, and radii
$L_i\downarrow0$ for which the rescaled varifolds converge to the
multiplicity-$Q$ plane. At these same scales there are nonzero measures
$\sigma_i$ carried by density-$Q$ points in a fixed compact ball, with
mass bounded above and away from zero and the estimate
\[
 \sigma_i(\mathbf B_r(x))\le Cr^s.
\]
Here $\mathbf B_r(x)$ is an ambient Euclidean ball. The limiting measure
of centers cannot be carried by an $(m-1)$-plane. This first reduction,
proved in Section~\ref{sec:flat-scales}, is the familiar geometric
starting point of dimension reduction; the challenge is to retain
these centers in a further, height-normalized limit.

The proof has four further stages. First, we construct smooth
reference graphs at the prescribed scales $L_i$. AE's minimal fitting
gives intervals of radii on which a centered height admits frequency
control. Section~\ref{sec:fitting} proves that these intervals cover
every sufficiently small prescribed flat scale, including their
transition regions. Section~\ref{sec:centered-blowup} then normalizes
the height by numbers $a_i\downarrow0$. The normalized second moment
is nonzero, the first moment tends to zero, and the reference graphs'
mean curvature and its first derivative are $o(a_i)$.
Both moment conclusions are needed: subtracting a reference graph
must neither remove the entire limit nor leave an arbitrary translation
in it.

Second, Section~\ref{sec:linearized} passes the first variations and
signed excess to measures on the limiting plane. The resulting data
consist of a nonnegative height measure, a positive-semidefinite tilt
measure, a mixed height--tilt measure, and a signed mass measure.
They satisfy linear first-variation identities, a positive-semidefinite
block inequality, and a one-sided comparison of mass and tilt. This
formulation avoids assuming that the limiting sheets can be globally
labeled or represented by a single-valued function.

Third, the density of the points carrying $\sigma_i$ must still be visible
after division by $a_i^2$. Ordinary uncorrected radial comparison
does not provide that precision on curved reference graphs.
Section~\ref{sec:density-centers} uses geodesic radial fields weighted by
the inverse reference-volume density. The resulting curvature error
can be averaged against $\sigma_i$ at every radius down to zero. For a
fixed $l>0$, its decisive integral is
\[
 \int_0^l t^{s-m}\,\dd t<\infty,
\]
which holds exactly when $s>m-1$. The centered first-moment estimate
then removes the translation terms and yields a density comparison
at every limiting center.

Finally, Section~\ref{sec:dimension-reduction} isolates an abstract
dimension-reduction statement for these measures. Its frequency is
monotone at the retained centers. A second blow-up makes that frequency
constant at enough independent centers. Equality in the
positive-semidefinite inequality then forces the height measure to
vanish, contrary to its normalization. This is the contradiction that
proves Theorem~\ref{thm:main}.

The prescribed-scale construction, the averaged curved radial
comparison, and the final measure-valued reduction are the parts needed
to strengthen AE's almost-everywhere conclusion to a dimension bound.
The center-manifold and frequency viewpoint belongs to the broader
regularity tradition described above; our specific local fitting and
signed-excess inputs are those of AE. The initial and final reductions
use the classical measure-theoretic framework of
\cite{Federer1969,Mattila1995}.

\subsection{Conventions}
All assertions are local in $U$, and supports are relative to $U$ until
a blow-up is taken. If $\eta_r(x)=x/r$, then $V_r=(\eta_r)_\#V$ is the
varifold pushforward, whose weight is $r^{-m}(\eta_r)_\#\mu$.
Translated blow-ups will be specified explicitly. Bold balls $\mathbf B_r$ are ambient, whereas
$B_r$ denotes a base-coordinate ball. The letters $C,c>0$ denote
constants depending only on previously fixed data and may change
between occurrences. Weak convergence of Radon measures is always
tested on continuous compactly supported functions. Positive
semidefiniteness of a symmetric matrix-valued measure means that its contraction
with every constant vector is a nonnegative measure.
